\section{Boolean ancestry, coverage, and collision}\label{sec:duality}

This section is valid for every positive loopless row-stochastic $P$;
reversibility is not used.

\subsection{The fair-geometric union dual}

At fitness two, \eqref{2.2} becomes $2x/(1+x)$.  This is a finite additive
graphical system in the sense of \citet{Harris1978}.  Let $L$ have the fair
geometric law
\[
 \Pp(L=\ell)=2^{-\ell},\qquad \ell\ge1,                       \tag{3.1}\label{3.1}
\]
and, conditionally on target $v$, draw $U_1,\ldots,U_L$
independently from row $P_v$.  For any mutant set $S$,
\[
 \Pp(U_1,\ldots,U_L\notin S)
 =\E(1-P_{vS})^L=\frac{1-P_{vS}}{1+P_{vS}}.                  \tag{3.2}\label{3.2}
\]
Thus the forward update at $v$ is exactly the additive Boolean map
\[
 x_v\longleftarrow x_{U_1}\vee\cdots\vee x_{U_L}.           \tag{3.3}\label{3.3}
\]
The transpose graphical map sends a nonempty ancestral set $A$ to
\[
 A\longmapsto
 \begin{cases}
 (A\setminus\{v\})\cup\{U_1,\ldots,U_L\},&v\in A,\\
 A,&v\notin A.
 \end{cases}                                                 \tag{3.4}\label{3.4}
\]
Because $P_{vv}=0$, the updated set in the first line omits $v$ and is
always nonempty; hence the dual remains a proper nonempty subset whenever it
starts there.

Using the same graphical marks forward and backward gives the exact duality
\[
 \Pp_S(X_t\cap A\ne\varnothing)
 =\Pp_A(S\cap A_t\ne\varnothing).                            \tag{3.5}\label{3.5}
\]

\begin{lemma}[ergodicity of the union dual]\label{lem:dual-ergodic}
If every off-diagonal entry of $P$ is positive, then the chain in
\eqref{3.4} is irreducible and aperiodic on the proper nonempty subsets of
$V$.
\end{lemma}

\begin{proof}
From any set $A$ of size at least two, choose a target $v\in A$, take $L=1$,
and sample $U_1\in A\setminus\{v\}$.  This removes $v$ and changes no other
element.  Repetition reaches a singleton.  Conversely, fix any proper
nonempty destination $B$ and choose $x\notin B$.  A singleton can first be
moved to $\{x\}$ using a one-sample update.  An update targeted at $x$ whose
finite geometric burst visits every element of $B$ and no element outside
$B$ then has output exactly $B$.  Every event just described has positive
probability because $P$ is positive off the diagonal and every finite value
of $L$ has positive probability.  Hence the chain is irreducible.  Finally,
every proper $A$ has a target outside $A$; choosing it leaves $A$ unchanged,
so every state has a positive self-loop and the chain is aperiodic.
\end{proof}

Let $\Pi_P$ denote its stationary law and put
\[
 m(P)=\E_{\Pi_P}|A|.                                         \tag{3.6}\label{3.6}
\]
For a set function $f$ and $i\notin S$, write
\[
 \Delta_i f(S)=f(S\cup\{i\})-f(S),\qquad
 \Delta_T=\prod_{i\in T}\Delta_i
 \quad(T\cap S=\varnothing).                                \tag{3.6a}\label{3.6a}
\]
The difference operators in the product commute.

\begin{proposition}[coverage representation]\label{prop:coverage}
The fixation committor at fitness two is
\[
 h_P(S)=\Pp_{A\sim\Pi_P}(A\cap S\ne\varnothing),            \tag{3.7}\label{3.7}
\]
and therefore
\[
 \rho_{\dB}(P,2)=\frac{m(P)}n.                              \tag{3.8}\label{3.8}
\]
Moreover, for every nonempty $T$ disjoint from $S$,
\[
 (-1)^{|T|+1}\Delta_T h_P(S)
 =\Pp_{\Pi_P}(A\cap S=\varnothing,\ T\subseteq A)\ge0.      \tag{3.9}\label{3.9}
\]
Thus $h_P$ is a normalized completely alternating coverage function.
\end{proposition}

\begin{proof}
The finite forward chain absorbs almost surely.  Indeed, from every
nonabsorbing mutant set, a sequence of resident-target deaths followed by
mutant-parent choices gives a positive-probability path to $V$; hence the
finite chain has no other closed communicating class.  Start the dual from $V$ in
\eqref{3.5}.  Its first active update sends $V$ to $V\setminus\{v\}$ because
all sampled parents differ from the target $v$; Lemma~\ref{lem:dual-ergodic}
then gives convergence to $\Pi_P$.  Letting $t\to\infty$, the left side
converges to the fixation probability from $S$, giving \eqref{3.7}.
Averaging $h_P(\{i\})=\Pp(i\in A)$ over $i$ proves
\eqref{3.8}.  For fixed $A$, the indicator
$\one_{\{A\cap S\ne\varnothing\}}$ has mixed difference
$(-1)^{|T|+1}\Delta_T$ equal to
$\one_{\{A\cap S=\varnothing,\ T\subseteq A\}}$; averaging proves
\eqref{3.9}.
\end{proof}

\subsection{A one-sample active chain}

The geometric burst can be factored into one-sample steps.  Two phase spaces
are required:
\[
 \mathcal Z_n=\{(C,v):C\subseteq V\setminus\{v\}\},\qquad
 \mathcal Y_n=\{(B,v):\varnothing\ne B\subseteq V\setminus\{v\}\}.
                                                                    \tag{3.10}\label{3.10}
\]
The pre-sampling space $\mathcal Z_n$ permits the empty cache created when a
singleton is stopped.  Let $A:\mathcal Z_n\to\mathcal Y_n$ adjoin one source
$i$ sampled from row $P_v$.  Let $R:\mathcal Y_n\to\mathcal Z_n$ be the fair
continue-or-retarget channel: on continue it sends $(B,v)$ to $(B,v)$, while
on stop it chooses $w$ uniformly from $B$ and sends $(B,v)$ to
$(B\setminus\{w\},w)$.  Row laws act on the right, so $K(P)=RA$ is a kernel
on $\mathcal Y_n$ and $M(P)=AR$ is a kernel on $\mathcal Z_n$.  Explicitly,
from $y=(B,v)$, with $b=|B|$, $K(P)$ is the following fair mixture:
\begin{enumerate}
\item sample $i$ from row $P_v$ and move to $(B\cup\{i\},v)$;
\item choose $w$ uniformly from $B$, sample $i$ from row $P_w$, and move to
      $((B\setminus\{w\})\cup\{i\},w)$.
\end{enumerate}
Each branch has probability $1/2$.  The first branch continues a geometric
burst; the second stops it and retargets the active ancestral mark.

\begin{lemma}[active collision identity]\label{lem:active}
The active kernel $K(P)$ is irreducible and aperiodic.  Let $\nu(P)$ be its
unique stationary law and put
\[
 H(B,v)=\frac1{|B|}.                                        \tag{3.11}\label{3.11}
\]
Then
\[
 \nu(P)H=\frac1{m(P)}=\frac1{n\rho_{\dB}(P,2)}.             \tag{3.12}\label{3.12}
\]
The map $P\mapsto K(P)$ is linear.
\end{lemma}

\begin{proof}
We include the marked-current calculation because it is the bridge from the
coverage law to the Hessian.  Let $T_v(A,B)$ be the transition kernel of one
union-dual update with target fixed at $v$.  Extend $\Pi_P$ by zero off the
proper nonempty subsets of $V$.  For $v\notin C$, including
$C=\varnothing$, define
\[
 \sigma_v(C)=\Pi_P(C\cup\{v\}),\qquad
 \eta_v(C)=\sum_A\Pi_P(A)T_v(A,C)-\Pi_P(C).                 \tag{3.13a}\label{3.13a}
\]
The subtracted term is exactly the passive current from $A=C$, for which
$v\notin A$ and no ancestral update occurs; hence $\eta_v(C)$ is the
nonnegative active current arriving at output $(C,v)$.  Stationarity of the
union dual, after summing its $n$ equally likely target updates, gives
\[
 \sum_{v\notin B}\eta_v(B)=|B|\Pi_P(B).                    \tag{3.13}\label{3.13}
\]
If $A_v$ adjoins one sample from row $P_v$, the fair-geometric resolvent
identity gives
\[
 \lambda_v=\frac{\sigma_v+\eta_v}{2},\qquad
 \eta_v=\lambda_vA_v.                                      \tag{3.14}\label{3.14}
\]
For the rectangular channels just defined, after the sample, continuing contributes
$\eta_v(B)/2$ to $(B,v)$.  Stopping, choosing $w$ uniformly from $B$, and
ending at $(B\setminus\{w\},w)$ contributes, by \eqref{3.13},
$\Pi_P(B)/2=\sigma_w(B\setminus\{w\})/2$.  Hence
$\lambda AR=\lambda$.  Thus $\lambda$ is stationary on $\mathcal Z_n$ for
$M=AR$.  Since $\eta=\lambda A$, associativity gives $\eta RA=\eta$; hence
$\eta$ is stationary on $\mathcal Y_n$ for the active kernel $K=RA$.  Both
$\sigma_v$ and $\eta_v$ have total mass
$\Pp_{\Pi_P}(v\in A)$, so
\[
 \sum_{v,C}\lambda_v(C)=m(P).                              \tag{3.15}\label{3.15}
\]
Normalizing gives the probability laws $\lambda/m(P)$ on $\mathcal Z_n$ and
$\nu(P)=\eta/m(P)$ on $\mathcal Y_n$.  The latter stationary law is unique:
if $|B|\ge2$, a stop at $w\in B$ followed by sampling a member of
$B\setminus\{w\}$ shrinks $B$ by one.  From a singleton state, one or two
positive stop moves reach any prescribed singleton--target pair $(\{j\},v)$
with $j\ne v$.  Positive continue moves then sample the remaining members of
any prescribed active set.  Hence $K(P)$ is irreducible; sampling an already
present source on a continue move gives a self-loop, so it is also aperiodic.
Dividing
\eqref{3.13} by $|B|$ and summing over $B$ yields
\[
 \nu(P)H=\sum_{k\ge1}\frac{k\Pi_P(|A|=k)}{m(P)}\frac1k
 =\frac1{m(P)}.                                             \tag{3.16}\label{3.16}
\]
Both active moves contain exactly one use of a row of $P$, proving linearity.
\end{proof}

At the complete kernel, put $N=n-1$.  Direct substitution gives
\[
 \nu_0(B,v)=\frac{|B|}{nN2^{N-1}},\qquad
 c_0:=\nu_0H=\frac{2^N-1}{N2^{N-1}}=\frac1{m_n}.             \tag{3.17}\label{3.17}
\]
This agrees with \eqref{2.4} at $r=2$.

The collision interpretation uses the conjugate sample-then-retarget chain,
not the retarget-then-sample phase order of $K=RA$.  Specifically,
\[
 \frac{\lambda}{m(P)}\xrightarrow{\ A\ }
 \frac{\eta}{m(P)}=\nu(P)\xrightarrow{\ R\ }
 \frac{\lambda}{m(P)},
 \qquad
 \frac{\eta}{m(P)}\xrightarrow{\ R\ }
 \frac{\lambda}{m(P)}\xrightarrow{\ A\ }
 \frac{\eta}{m(P)}.                                      \tag{3.18a}\label{3.18a}
\]
Thus $M(P)=AR$ has stationary law $\lambda/m(P)$, while $K(P)=RA$ has
stationary law $\eta/m(P)=\nu(P)$.  In one $M(P)$ step, let $I$ be the source
drawn by $A$.  The post-sampling set $B$ contains $I$ and has law $\nu(P)$.
If the subsequent fair coin in $R$ stops, it chooses the new target $W$
uniformly from that same set $B$.  Consequently
\[
 \Pp(W=I\mid B,\mathrm{stop})=\frac1{|B|},
\]
and averaging with \eqref{3.12} gives
\[
 \Pp(\mathrm{stop},W=I)=\frac1{2m(P)},\qquad
 \Pp(W=I\mid\mathrm{stop})=\frac1{m(P)}.                  \tag{3.18}\label{3.18}
\]
Thus maximizing fixation is equivalent to minimizing this stationary
collision rate.
